\documentclass[../DoAn.tex]{subfiles}
\begin{document}

Phụ lục này mô tả bộ kịch bản kiểm thử được dùng trong Chương 5. Các kịch bản được nhóm theo yếu tố cần kiểm tra để tránh trộn lẫn ảnh hưởng của vị trí căn cứ, năng lượng, vật cản, địa hình và giả định chi phí quay.

\section{Tổng quan các nhóm kịch bản}

\begin{table}[H]
\centering
\caption{Danh mục nhóm kịch bản kiểm thử}
\small
\setlength{\tabcolsep}{4pt}
\renewcommand{\arraystretch}{1.08}
\begin{tabular}{|C{0.10\textwidth}|L{0.22\textwidth}|L{0.58\textwidth}|}
\hline
\textbf{Nhóm} & \textbf{Kịch bản} & \textbf{Mục tiêu kiểm tra} \\
\hline
A & A01--A10 & Ảnh hưởng của vị trí căn cứ trên cùng một bản đồ chính. \\
\hline
B & B01--B09 & Ảnh hưởng của dung lượng năng lượng tối đa. \\
\hline
C & C01--C10 & Ảnh hưởng của mật độ vật cản và tính liên thông của bản đồ. \\
\hline
D & D13--D14 & Ảnh hưởng của vật cản động tới đường đi hoặc đường quay về. \\
\hline
E & E01--E10 & Ảnh hưởng của chi phí địa hình có trọng số. \\
\hline
F & F01--F05 & Khả năng duy trì quy trình nhiệm vụ khi kích thước bản đồ tăng. \\
\hline
G & G01--G11 & Ảnh hưởng của cấu trúc phân bố vật cản khi mật độ gần như giữ nguyên. \\
\hline
H & H01--H06, H17\_D14 & Đối chứng toàn hướng, trong đó chi phí quay được bỏ qua ở mức lập kế hoạch. \\
\hline
\end{tabular}
\end{table}

\section{Nhóm A, B và C}

Nhóm A dùng cùng một bản đồ chính kích thước 35$\times$30 và thay đổi vị trí căn cứ. Nhóm B giữ bản đồ nhỏ ổn định nhưng giảm dần năng lượng tối đa. Nhóm C thay đổi mật độ vật cản và phân biệt bản đồ còn liên thông với bản đồ bị chia cắt.

\begin{table}[H]
\centering
\caption{Tóm tắt các nhóm A, B và C}
\small
\setlength{\tabcolsep}{4pt}
\renewcommand{\arraystretch}{1.08}
\begin{tabular}{|C{0.10\textwidth}|L{0.25\textwidth}|L{0.55\textwidth}|}
\hline
\textbf{Nhóm} & \textbf{Biến thay đổi} & \textbf{Ghi chú} \\
\hline
A & Vị trí căn cứ & Các vị trí gồm trung tâm, góc, biên và lệch tâm. \\
\hline
B & Năng lượng tối đa & Năng lượng giảm từ 260 xuống 50 để quan sát số lần quay về và sạc lại. \\
\hline
C & Mật độ vật cản & Có cả bản đồ liên thông và bản đồ có vùng không thể tiếp cận. \\
\hline
\end{tabular}
\end{table}

\section{Nhóm D, E, F và G}

Nhóm D đưa vật cản động vào mô phỏng. Nhóm E thay đổi chi phí địa hình. Nhóm F tăng kích thước bản đồ. Nhóm G giữ mật độ vật cản gần như tương đương nhưng thay đổi cấu trúc hình học của vật cản.

\begin{table}[H]
\centering
\caption{Tóm tắt các nhóm D, E, F và G}
\small
\setlength{\tabcolsep}{4pt}
\renewcommand{\arraystretch}{1.08}
\begin{tabular}{|C{0.10\textwidth}|L{0.25\textwidth}|L{0.55\textwidth}|}
\hline
\textbf{Nhóm} & \textbf{Biến thay đổi} & \textbf{Ghi chú} \\
\hline
D & Vật cản động & Kiểm tra tình huống đường đi hoặc đường quay về bị ảnh hưởng. \\
\hline
E & Trọng số địa hình & Kiểm tra khi các ô hợp lệ có chi phí khác nhau. \\
\hline
F & Kích thước bản đồ & Kiểm tra quy trình nhiệm vụ trên bản đồ lớn hơn. \\
\hline
G & Cấu trúc vật cản & Kiểm tra sự khác biệt giữa các bản đồ có mật độ gần tương đương nhưng hình học khác nhau. \\
\hline
\end{tabular}
\end{table}

\section{Nhóm H: đối chứng toàn hướng}

Nhóm H dùng mô hình toàn hướng lý tưởng, trong đó chi phí quay được bỏ qua ở mức lập kế hoạch. Nhóm này không thay thế mô hình chính mà đóng vai trò đối chứng, giúp đánh giá ảnh hưởng của chi phí quay tới số bước, năng lượng, số lần quay về và trạng thái kết thúc nhiệm vụ.

\end{document}
